\section{Cancelaciones locales con un exterior común}
Especificamos una implicación formal para aplicar la parte 05. El exterior se obtiene
quitando todos los sitios marcados antes de elegir las restricciones locales activas.
Sean $\zeta_i$ campos independientes de esa elección. Una cara es una elección de
restricciones. Los coeficientes exteriores congelados deben ser los mismos en todas las caras.
El campo afín de orden cero se absorbe en la ley basal de cada componente.
A costes $l\ge1$ supongamos la descomposición de \emph{Patrones finitos}, parte 03.
Cada monomio sólo depende de los grupos que se le asignan. La acción ordinaria tiene grado
$\le l+1$, el prefactor ordinario grado $\le l$. La excepción de dos fibras tiene una
amplitud común en acción y prefactor como en la parte 05 y soporte $\le l$.
Las demás excepciones tienen soporte $\le l-1$. En soporte máximo, los multiplicadores
de la excepción de dos fibras son lineales en grupos privados distintos.
\paragraph{Bosque, a campo fijo.}
Cada grupo forestal tiene partición de Bethe exacta y las mismas respuestas de campo.
En la silla inclinada $m$, $Z=X-m$ tiene covarianza por bloques $V$.
Tras sacar los términos afines, el logaritmo residual a campo fijo tiene la forma
\[
 \log\E\exp\{R(m+Z)-R(m)-R'(m)Z+A(m+Z)-A(m)\}
   +\tfrac12\log\det(I-VR''(m)).
\]
En esta expresión se considera el sector ordinario. Sean $p$ el número de patas centradas,
$q_0$ el número de factores de media y $b$ el número de bloques de cumulantes locales.
Sean $N_R,N_A$ los números de vértices de acción y prefactor, y $w$ el coste total.
En un diagrama conexo, con exceso de ciclos $L=p-(N_R+N_A+b)+1$, el número $d$
de grupos distintos satisface
\[
       p+q_0\le w+N_R,\qquad d\le q_0+b\le w+1-N_A-L .
\]
Cada vértice de acción y cada bloque tienen al menos dos patas. Cada vértice de
prefactor tiene al menos una. Con $N_A=0$ o $1$ no puede haber un árbol, pues
tendría al menos dos hojas. El diagrama vacío tiene coeficiente cero en el logaritmo.
Para $N_A=0,L=1$ todos los grados valen dos. Los ciclos con $k$ bloques aportan
$\operatorname{tr}[(VR''(m))^k]/(2k)$: la división cuenta las rotaciones y la
reversión, válida también como contracción gaussiana para $k=1,2$.
El medio logaritmo del determinante cancela la suma. Los términos ordinarios
que sobreviven tienen soporte $d\le w-1$.
Al reponer la silla, la derivada por el grupo $g$ de $R_l$ añade a lo sumo $l$
grupos además de $g$. Las derivadas locales de la media siguen en $g$. Por inducción
en el coste, la reposición a coste $a$ sólo añade hasta $a$ grupos. Se conserva $d\le w-1$.
\paragraph{Raíz de dos fibras.}
Para un coeficiente conexo cuyo grafo de incidencias tiene exceso $L_I$, con $h_b$
grupos distintos en la interacción $b$, vale $d=1+w-\sum_b(l_b+1-h_b)-L_I$.
Con una raíz excepcional, el único caso que puede alcanzar $d=w$ es un árbol
con las otras interacciones multilineales de soporte máximo. Tomemos una interacción hoja
cuyos grupos, salvo el de enlace, son privados. Al anular sólo su amplitud, dejando
activas las del resto, los grupos privados y el componente restante se separan.
En el valor exacto, derivar en esa amplitud inserta la ocupación del enlace por las
medias privadas. En la silla, la variación en el grupo de enlace es la de su campo por
el mismo producto. La Hessiana explícita de la hoja sólo cruza bloques distintos
en ese punto y su traza contra la inversa por bloques es cero. Se deriva también
el prefactor y el determinante del componente restante al variar el campo.
Si los factores privados responden, la derivada de su valor estacionario ya está
incluida por estacionariedad y su prefactor forestal local es exacto.
Esta identidad con las demás amplitudes libres permite derivarlas y podar por inducción.
Queda la raíz de la parte 05, cuya diferencia y derivadas de campo se anulan.
Dos déficits o una interacción de soporte $\le l-1$ ya dan $d\le w-1$.
\paragraph{Un único ciclo.}
Con $c$ grupos que deben aparecer, en grado $w=c-1>0$ sólo cabe un árbol
de interacciones de soporte máximo. El caso $c=1,w=0$ se lee directamente en el grupo.
En ese árbol la partición en cumulantes conectados no puede separar las incidencias
de un grupo: hacerlo desconectaría el árbol. Sus ramas son derivadas de campo locales.
Si el grupo que queda tiene un ciclo binario $C$, la transición de $i$ a $i+1$ tiene
filas $(1-t_i,t_i),(1,0)$ con $t_i=x_{i+1}/(1-x_i)$. La traza del producto
vale $1+\theta$ con $\theta=\prod_{i\in C}(-t_i)$. En la corrección gaussiana,
las regresiones tienen esas mismas pendientes negativas y cerrar el sistema de
innovaciones da jacobiano $1-\theta$. Los factores de rama se cancelan;
la diferencia local es $\log(1-\theta^2)$ en la rama del origen.
Con $x_i=\zeta_i b_i$ holomorfo, conserva $\prod_{i\in C}\zeta_i^2$ incluso tras
derivar con $\zeta_i\partial_{\zeta_i}$. Si una cara rompe el ciclo, queda un bosque.
\paragraph{Hipótesis de uso.}
Al poner $\zeta_i=0$ se exige además que las dos caras de ese sitio coincidan
en el normalizador completo. La proyección en sitios ausentes da entonces un factor
$\zeta_i$ por sitio, además de la potencia extra ya identificada en el ciclo presente.
Localmente, para exclusión de capacidad uno, las matrices de covarianza normalizadas
tienen diagonal uno y entradas cruzadas $-x_j/(1-x_i)$ en cada factor.
Para $\max_f\sum_{i\in f}|\zeta_i e^{h_i}|$ pequeño y grado de sitio acotado,
las ecuaciones de mensajes son contractivas alrededor de uno y estos bloques
son próximos a la identidad. Esto justifica la rama local cerca de campos cero.
El uso en el predictor físico exige cotejar el factor exterior y su normalización
en todas las caras. La parte 07 comprueba esa identidad y la parte 08 explicita
qué radios se usan para estimar un soporte cuyo número de marcas aumenta.
